\documentclass[11pt,a4paper]{article}
\usepackage[utf8]{inputenc}
\usepackage[T1]{fontenc}
\usepackage[french]{babel}
\usepackage{amsmath,amssymb,amsfonts}
\usepackage{geometry}
\usepackage{tikz}
\usepackage{fancyhdr}

\geometry{margin=2.5cm}

\pagestyle{fancy}
\fancyhf{}
\fancyfoot[C]{\thepage}

\begin{document}

\noindent
\textbf{D.BENMAKHLOUF} \\
Classe : Seconde \\
Année scolaire : 2025-2026

\vspace{0.5cm}

\begin{center}
\textbf{\Large Évaluation : Forme canonique et factorisation du second degré}
\end{center}

\vspace{0.3cm}

\section*{Sujet A}

On considère la fonction polynôme $f$ définie sur $\mathbb{R}$ par :
$$f(x) = x^2 - 2x - 3$$

\begin{enumerate}
    \item \textbf{Forme canonique et sommet}
    \begin{enumerate}
        \item Vérifier que pour tout réel $x$, on a : $f(x) = (x - 1)^2 - 4$.
        \item En déduire les coordonnées du sommet $S$ de la parabole représentant $f$.
        \item Quel est le sens de variation de la fonction $f$ sur $\mathbb{R}$ ? Dresser son tableau de variations.
    \end{enumerate}
    
    \item \textbf{Factorisation et racines}
    \begin{enumerate}
        \item Montrer que $f(x) = (x + 1)(x - 3)$ pour tout réel $x$.
        \item Résoudre dans $\mathbb{R}$ l'équation $f(x) = 0$. En déduire les intersections de la parabole avec l'axe des abscisses.
        \item Déterminer les coordonnées du point d'intersection de la parabole avec l'axe des ordonnées.
    \end{enumerate}
    
    \item \textbf{Lecture graphique et interprétation}
    \begin{quotation}
    Dans un repère orthonormé $(O, \vec{i}, \vec{j})$, la courbe $\mathcal{C}_f$ est tracée ci-dessous. Retrouver graphiquement le sommet $S$ et les racines de $f$.
    \end{quotation}
\end{enumerate}

\begin{center}
\begin{tikzpicture}[scale=0.85]
    \draw[very thin, lightgray!60] (-3, -5) grid (5, 5);
    \draw[->, thick] (-3.2, 0) -- (5.2, 0) node[right] {$x$};
    \draw[->, thick] (0, -5.2) -- (0, 5.2) node[above] {$y$};
    \node at (0,0) [below left] {$O$};
    \node at (1,0) [below] {$1$};
    \node at (0,1) [left] {$1$};

    \begin{scope}\clip (-3.0,-5.0) rectangle (5.0,5.0);
    \draw[domain=-1.3:3.3, samples=100, smooth, blue, very thick] plot (\x, {(\x)*(\x) - 2*(\x) - 3});
    \end{scope}
    \node[blue, above right] at (2, 2) {$\mathcal{C}_f$};

    \fill[red] (1, -4) circle (2.5pt);
    \node[red, below right] at (1, -4) {$S(1; -4)$};

% [autofix neutralisé]     \fill[darkgreen!60!black] (-1, 0) circle (2.0pt);
% [autofix neutralisé]     \node[darkgreen!60!black, above left] at (-1, 0) {$(-1; 0)$};

% [autofix neutralisé]     \fill[darkgreen!60!black] (3, 0) circle (2.0pt);
% [autofix neutralisé]     \node[darkgreen!60!black, above right] at (3, 0) {$(3; 0)$};

    \fill[orange!80!black] (0, -3) circle (2.0pt);
    \node[orange!80!black, left] at (0, -3) {$(0; -3)$};
\end{tikzpicture}
\end{center}

\newpage

\section*{Corrigé détaillé du Sujet A}

\begin{enumerate}
    \item \textbf{Forme canonique et sommet}
    \begin{enumerate}
        \item Développons l'expression proposée :
        $$(x - 1)^2 - 4 = (x^2 - 2x + 1) - 4 = x^2 - 2x - 3 = f(x)$$
        L'égalité est donc vérifiée.
        \item La forme canonique est $f(x) = a(x - \alpha)^2 + \beta$ avec $a = 1$, $\alpha = 1$ et $\beta = -4$. Le sommet $S$ de la parabole a pour coordonnées $(\alpha ; \beta)$, soit $S(1 ; -4)$.
        \item Comme le coefficient $a = 1 > 0$, la parabole est tournée vers le haut. La fonction est donc décroissante sur l'intervalle $]-\infty ; 1]$ puis croissante sur $[1 ; +\infty[$.
        
        \vspace{0.2cm}
        \noindent\textbf{Tableau de variations :} \\
        \begin{center}
        \begin{tabular}{|c|ccccc|}
        \hline
        $x$ & $-\infty$ & & $1$ & & $+\infty$ \\
        \hline
        & $+\infty$ & & & & $+\infty$ \\
        $f(x)$ & & $\searrow$ & & $\swarrow$ & \\
        & & & $-4$ & & \\
        \hline
        \end{tabular}
        \end{center}
    \end{enumerate}
    
    \item \textbf{Factorisation et racines}
    \begin{enumerate}
        \item Développons la forme factorisée :
        $$(x + 1)(x - 3) = x^2 - 3x + x - 3 = x^2 - 2x - 3 = f(x)$$
        L'égalité est établie.
        \item Résolvons $f(x) = 0$ à l'aide de la forme factorisée :
        $$(x + 1)(x - 3) = 0$$
        Un produit de facteurs est nul si et seulement si l'un de ses facteurs est nul :
        $$x + 1 = 0 \quad \text{ou} \quad x - 3 = 0$$
        C'est-à-dire $x = -1$ ou $x = 3$. Les solutions sont $-1$ et $3$, ce qui correspond aux points d'intersection de la courbe avec l'axe des abscisses : $(-1 ; 0)$ et $(3 ; 0)$.
        \item Pour l'intersection avec l'axe des ordonnées, on calcule $f(0)$ :
        $$f(0) = 0^2 - 2(0) - 3 = -3$$
        Le point d'intersection a pour coordonnées $(0 ; -3)$.
    \end{enumerate}
    
    \item \textbf{Lecture graphique} \\
    Le point le plus bas de la parabole se lit directement en $S(1 ; -4)$. Les points d'intersection avec l'axe des abscisses se lisent aux abscisses $-1$ et $3$.
\end{enumerate}

\vspace{0.5cm}

\section*{Sujet B (Variante d'entraînement)}

On considère la fonction polynôme $g$ définie sur $\mathbb{R}$ par :
$$g(x) = x^2 - 4x + 3$$

\begin{enumerate}
    \item \textbf{Forme canonique et extremum}
    \begin{enumerate}
        \item Démontrer que pour tout réel $x$, $g(x) = (x - 2)^2 - 1$.
        \item Préciser les coordonnées du sommet de la parabole représentant $g$ et donner son extremum.
        \item Dresser le tableau de variations complet de la fonction $g$.
    \end{enumerate}
    \item \textbf{Factorisation et intersections}
    \begin{enumerate}
        \item Vérifier que $g(x) = (x - 1)(x - 3)$ pour tout réel $x$.
        \item Déterminer les antécédents de $0$ par la fonction $g$ et en déduire les points d'intersection de la courbe avec l'axe $(O\vec{x})$.
        \item Calculer l'ordonnée à l'origine de la parabole.
    \end{enumerate}
\end{enumerate}

\end{document}